\clearpage\section*{\centering REFERÊNCIAS}\addcontentsline{toc}{section}{REFERÊNCIAS}
\begin{singlespace}\setlength{\parindent}{0pt}\setlength{\parskip}{\baselineskip}\raggedright
ASSOCIAÇÃO BRASILEIRA DE NORMAS TÉCNICAS. \textbf{ABNT NBR 6034: informação e documentação: índice: apresentação}. Rio de Janeiro: ABNT, 2004a.

ASSOCIAÇÃO BRASILEIRA DE NORMAS TÉCNICAS. \textbf{ABNT NBR 12225: informação e documentação: lombada: apresentação}. Rio de Janeiro: ABNT, 2004b.

ASSOCIAÇÃO BRASILEIRA DE NORMAS TÉCNICAS. \textbf{ABNT NBR 6024: informação e documentação: numeração progressiva das seções de um documento: apresentação}. Rio de Janeiro: ABNT, 2012a.

ASSOCIAÇÃO BRASILEIRA DE NORMAS TÉCNICAS. \textbf{ABNT NBR 6027: informação e documentação: sumário: apresentação}. Rio de Janeiro: ABNT, 2012b.

ASSOCIAÇÃO BRASILEIRA DE NORMAS TÉCNICAS. \textbf{ABNT NBR 6022: informação e documentação: artigo em publicação periódica técnica e/ou científica: apresentação}. Rio de Janeiro: ABNT, 2018a. Disponível em: \url{https://www2.ufrb.edu.br/bcet/images/NBR_6022-2018.pdf}. Acesso em: 8 out. 2026.

ASSOCIAÇÃO BRASILEIRA DE NORMAS TÉCNICAS. \textbf{ABNT NBR 6023: informação e documentação: referências: elaboração}. Rio de Janeiro: ABNT, 2018b. Versão corrigida 2, de 24 set. 2020.

ASSOCIAÇÃO BRASILEIRA DE NORMAS TÉCNICAS. \textbf{ABNT NBR 6028: informação e documentação: resumo, resenha e recensão: apresentação}. Rio de Janeiro: ABNT, 2021.

ASSOCIAÇÃO BRASILEIRA DE NORMAS TÉCNICAS. \textbf{ABNT NBR 10520: informação e documentação: citações em documentos: apresentação}. Rio de Janeiro: ABNT, 2023.

ASSOCIAÇÃO BRASILEIRA DE NORMAS TÉCNICAS. \textbf{ABNT NBR 14724: informação e documentação: trabalhos acadêmicos: apresentação}. Rio de Janeiro: ABNT, 2024.

BRASIL. Conselho Nacional de Saúde. \textbf{Resolução nº 510, de 7 de abril de 2016}. Brasília, DF: CNS, 2016. Disponível em: \url{https://bvsms.saude.gov.br/bvs/saudelegis/cns/2016/res0510_07_04_2016.html}. Acesso em: 8 out. 2026.

INSTITUTO BRASILEIRO DE GEOGRAFIA E ESTATÍSTICA. \textbf{Normas de apresentação tabular}. 3. ed. Rio de Janeiro: IBGE, 1993. Disponível em: \url{https://biblioteca.ibge.gov.br/visualizacao/livros/liv23907.pdf}. Acesso em: 8 out. 2026.

PRATHER, James et al. The widening gap: the benefits and harms of generative AI for novice programmers. In: INTERNATIONAL COMPUTING EDUCATION RESEARCH CONFERENCE, 2024. \textbf{Proceedings [...]}. ACM, 2024. DOI: \url{https://doi.org/10.1145/3632620.3671116}. Versão dos autores: \url{https://arxiv.org/abs/2405.17739}.

QUEIROZ, Francisca Karolina; LIMA, Márcia Sampaio. Uso do ChatGPT na priorização de requisitos: uma experiência educacional em engenharia de software. In: SIMPÓSIO BRASILEIRO DE EDUCAÇÃO EM COMPUTAÇÃO, 2025. \textbf{Anais [...]}. Sociedade Brasileira de Computação, 2025. DOI: \url{https://doi.org/10.5753/educomp.2025.5370}.

UNIVERSIDADE FEDERAL DO ESPÍRITO SANTO. Sistema Integrado de Bibliotecas. \textbf{Normalização}. Atualizado em 10 jul. 2026. Disponível em: \url{https://biblioteca.ufes.br/normalizacao}. Acesso em: 8 out. 2026.
\end{singlespace}
